\documentclass[aspectratio=169,10pt]{beamer}
\usetheme{Madrid}
\usecolortheme{default}

\usepackage{hyperref}
\usepackage{tikz}
\usepackage{amsmath,amssymb}

% --- Custom Colors ---
\definecolor{unibblue}{RGB}{0,51,102}
\definecolor{unibgold}{RGB}{212,175,55}
\definecolor{unibgreen}{RGB}{0,128,0}

\setbeamercolor{palette primary}{bg=unibblue,fg=white}
\setbeamercolor{palette secondary}{bg=unibblue!85,fg=white}
\setbeamercolor{palette tertiary}{bg=unibblue!70,fg=white}
\setbeamercolor{structure}{fg=unibblue}
\setbeamercolor{block title}{bg=unibblue!85,fg=white}
\setbeamercolor{block body}{bg=unibblue!10,fg=black}
\setbeamercolor{alertblock title}{bg=unibgold!90,fg=black}
\setbeamercolor{alertblock body}{bg=unibgold!15,fg=black}

\title{Persamaan Diferensial}
\subtitle{Minggu 10: Metode Pemisahan Variabel (*Separation of Variables*) pada PDP}
\author{Novalio Daratha \& Muhammad Arfan}
\institute{Program Studi Teknik Elektro\\Universitas Bengkulu}
\date{2026}

\begin{document}

\begin{frame}
  \titlepage
\end{frame}

\begin{frame}{Capaian Pembelajaran (CPMK 4)}
  \begin{itemize}
    \item \textbf{CPMK 4.1:} Memahami prinsip dasar pemisahan variabel: memecah PDP multi-variabel menjadi 2 atau lebih PDB independen.
    \item \textbf{CPMK 4.2:} Memahami peran konstanta pemisahan (*separation constant* $k$ atau $-\lambda^2$) dan justifikasi fisiknya.
    \item \textbf{CPMK 4.3:} Menerapkan syarat batas (*Boundary Conditions*) Dirichlet dan Neumann pada PDB spasial (Masalah Nilai Eigen).
    \item \textbf{CPMK 4.4:} Menggabungkan solusi eigen menggunakan prinsip superposisi dan koefisien Fourier.
  \end{itemize}
\end{frame}

\begin{frame}{Daftar Isi}
  \tableofcontents
\end{frame}

\section{Prinsip Dasar Pemisahan Variabel}
\begin{frame}{Konsep Pemisahan Produk Fungsi}
  Misalkan kita ingin mencari solusi $u(x,t)$ dari PDP linear:
  \begin{block}{Asumsi Pemisahan Produk}
    $$ {\color{unibblue} u(x,t) = X(x) \cdot T(t)} $$
    Fungsi multi-variabel $u(x,t)$ diasumsikan merupakan perkalian antara fungsi murni spasial $X(x)$ dan fungsi murni temporal $T(t)$.
  \end{block}
  \pause
  Turunan parsial menjadi turunan total biasa:
  $$ \frac{\partial u}{\partial t} = X(x) \frac{dT}{dt} = X T', \quad \frac{\partial^2 u}{\partial x^2} = \frac{d^2X}{dx^2} T(t) = X'' T $$
\end{frame}

\section{Studi Kasus: Pemisahan pada Persamaan Panas 1D}
\begin{frame}{Penurunan Dua PDB Terpisah}
  Tinjau: $\frac{\partial u}{\partial t} = \alpha \frac{\partial^2 u}{\partial x^2}$.
  \pause
  Substitusi $u(x,t) = X(x)T(t)$:
  $$ X T' = \alpha X'' T \implies \frac{T'}{\alpha T} = \frac{X''}{X} $$
  \pause
  \begin{alertblock}{Argumen Konstanta Pemisahan}
    Ruas kiri hanya bergantung pada $t$, ruas kanan hanya bergantung pada $x$. Agar persamaan ini bernilai sama untuk setiap $x$ dan $t$ sembarang, kedua ruas harus bernilai konstan yang sama ($-\lambda^2$):
    $$ \frac{X''}{X} = -\lambda^2 \implies {\color{unibblue} X'' + \lambda^2 X = 0} \quad \text{(PDB Spasial)} $$
    $$ \frac{T'}{\alpha T} = -\lambda^2 \implies {\color{unibblue} T' + \alpha\lambda^2 T = 0} \quad \text{(PDB Waktu)} $$
  \end{alertblock}
\end{frame}

\begin{frame}{Penyelesaian Masalah Nilai Eigen Spasial ($X(x)$)}
  Untuk konduktor panjang $L$ dengan ujung digroundkan $u(0,t) = 0$ dan $u(L,t) = 0$:
  $$ X(0) = 0, \quad X(L) = 0 $$
  \pause
  Solusi umum PDB spasial $X'' + \lambda^2 X = 0$:
  $$ X(x) = C_1 \cos(\lambda x) + C_2 \sin(\lambda x) $$
  \pause
  \begin{itemize}
    \item $X(0) = 0 \implies C_1 \cos(0) + C_2 \sin(0) = 0 \implies C_1 = 0$. \pause
    \item $X(L) = 0 \implies C_2 \sin(\lambda L) = 0$. Agar solusi non-trivial ($C_2 \neq 0$):
    $$ \sin(\lambda L) = 0 \implies \lambda L = n\pi \implies {\color{unibblue} \lambda_n = \frac{n\pi}{L}}, \quad n = 1, 2, 3, \dots $$
  \end{itemize}
  \pause
  \begin{block}{Fungsi Eigen (*Eigenfunctions*)}
    $$ X_n(x) = \sin\left(\frac{n\pi x}{L}\right) $$
  \end{block}
\end{frame}

\begin{frame}{Penyelesaian Bagian Waktu ($T_n(t)$) \& Superposisi}
  PDB Waktu: $T_n' + \alpha \left(\frac{n\pi}{L}\right)^2 T_n = 0$.
  \pause
  Solusi eksponensial meluruh:
  $$ T_n(t) = B_n e^{-\alpha \left(\frac{n\pi}{L}\right)^2 t} $$
  \pause
  \begin{alertblock}{Solusi Total Superposisi}
    $$ {\color{unibblue} u(x,t) = \sum_{n=1}^{\infty} B_n \sin\left(\frac{n\pi x}{L}\right) e^{-\alpha \left(\frac{n\pi}{L}\right)^2 t}} $$
    dengan $B_n = \frac{2}{L} \int_{0}^{L} f(x) \sin\left(\frac{n\pi x}{L}\right) dx$ diperoleh dari syarat awal $u(x,0) = f(x)$.
  \end{alertblock}
\end{frame}

\section{Kuis \& Rangkuman}
\begin{frame}{Kuis Interaktif 10}
  \textbf{Pertanyaan:} Mengapa konstanta pemisahan harus dipilih bernilai negatif ($-\lambda^2$) dan bukan positif ($+\lambda^2$)?
  \pause
  \begin{block}{Jawaban:}
    Secara fisik:
    \begin{enumerate}
      \item Jika konstanta positif ($+\lambda^2$), solusi spasial menjadi fungsi hiperbolik $\sinh(\lambda x), \cosh(\lambda x)$ yang tidak dapat memenuhi syarat batas nol di kedua ujung.
      \item Solusi waktu akan menjadi $T(t) \propto e^{+\lambda^2 \alpha t}$, yang berarti suhu atau energi membesar tanpa batas menuju tak terhingga (tidak stabil dan melanggar hukum termodinamika/konservasi energi).
    \end{enumerate}
  \end{block}
\end{frame}

\begin{frame}{Rangkuman Materi Minggu 10}
  \begin{enumerate}
    \item Pemisahan variabel mentransformasikan 1 PDP menjadi kumpulan PDB yang dapat diselesaikan secara independen.
    \item Syarat batas menghasilkan nilai eigen diskrit ($\lambda_n$) dan fungsi harmonik eigen ($X_n(x)$).
    \item Superposisi tak hingga menyatukan seluruh modus gelombang untuk mencocokkan kondisi awal sembarang.
  \end{enumerate}
\end{frame}

\end{document}
